% =====================================================================
% III. RESULTS AND DISCUSSION  (parts a-i)
% Numbers are taken from results/part_a.txt ... results/part_i.txt.
% LLM declaration: text level 3; see the appendix note.
% =====================================================================

\section{Results and discussion}\label{sec:results}

\subsection{Ordinary least squares}\label{sec:res_ols}

\subsubsection{Error against polynomial degree}

Figure~\ref{fig:a_mse_r2} shows MSE and $R^2$ against degree for $n = 100$ and $\sigma = 0.1$. The median training MSE falls steadily and drops below $\sigma^2 = 0.01$ from degree 11. This is expected: the fit absorbs part of the noise, and for a correct model the expected training MSE is $\sigma^2(1 - p/n)$ while the expected test MSE is $\sigma^2(1 + p/n)$ \cite[Sec.~3.7]{HjorthJensen2026}. The median test MSE never reaches $\sigma^2$. It flattens from about degree 10, between 0.0125 and 0.0142, and the lowest median, 0.0125 at degree 14, lies well inside the interquartile band of the neighbouring degrees. Between degree 1 and 15 with $n = 100$ there is little overfitting; it becomes clear for smaller $n$ or higher degree (Fig.~\ref{fig:a_n_sigma}).

The MSE falls mainly from odd to even degree (median test MSE 0.0453 and 0.0455 at degrees 2 and 3, 0.0265 and 0.0276 at degrees 4 and 5). Runge's function is even. The even powers can describe it, while the odd powers are odd functions and can only fit asymmetry, which here comes from the noise and the uneven placement of the $x$ values. An odd power therefore gives no gain on the test data, and a small loss.

\begin{figure*}
    \centering
    \includegraphics[width=0.85\textwidth]{a_mse_r2_degree.pdf}
    \caption{MSE (left, logarithmic) and $R^2$ (right) against polynomial degree for OLS, $n = 100$, $\sigma = 0.1$. Solid lines: median over 200 repetitions with new data and a new split each time; bands: interquartile range. Dashed lines: the single split with seed 42. Dotted line: $\sigma^2$.}
    \label{fig:a_mse_r2}
\end{figure*}

\subsubsection{How representative is one split?}

The single split (seed 42) lies below the interquartile band, and at degree 12 its test MSE, 0.0086, is even below $\sigma^2$. A test MSE below $\sigma^2$ cannot be caused by overfitting. The 20 test points happened to get less noise than average (mean squared noise $7.8\times10^{-3}$), and none of them lies at $|x| < 0.2$, around the peak. Since $f$ is known, we can compute the MSE the fitted model would get on new data: 0.0140 at degree 10, against 0.0094 on the test set. The same split gives $R^2 < 0$ up to degree 5, although its MSE there is at or below the median. The variance of $y$ in the test set is only 0.023, against 0.103 in the training set, and the same MSE divided by a four times smaller spread gives a much lower $R^2$. $R^2$ has randomness in both numerator and denominator, and its baseline is the mean of the test set itself, which a real model could not know. One small split thus makes the MSE too optimistic and $R^2$ too pessimistic, which is why we repeat the experiments, and later use resampling.

\subsubsection{Extreme repetitions and extrapolation}\label{sec:res_extreme}

At degree 15 the median test MSE is 0.013, while the mean is 0.106. Our hypothesis was that the extreme repetitions are those where a test point lies outside the interval covered by the training points, so that the high-degree polynomial must extrapolate. In 68 of 200 repetitions at least one test point lay outside, typically only a little (median distance 0.014). From degree 10, all ten worst repetitions extrapolate, and without the test points outside the training interval the mean at degree 15 drops to 0.015 (Table~\ref{tab:a_extrapolation}; see also Fig.~\ref{fig:a_extrapolation} in the appendix). At degree 20 extrapolation does not explain everything: the worst remaining repetitions have their worst test point in a gap of 0.08--0.14 in the training data, against a median spacing of 0.017, all at $|x| > 0.8$, where the prediction misses $f$ by up to 5. Inside a gap near the edge the polynomial is held from one side only and the high powers are largest. The large errors are caused by where the training points fell, not by the noise: at high degree the fit depends strongly on the particular sample.

\begin{table}
    \centering
    \caption{Test MSE over 200 repetitions ($n = 100$, $\sigma = 0.1$), with all test points and with only those inside the interval covered by the training data.}
    \label{tab:a_extrapolation}
    \begin{tabular}{rcccc}
        \toprule
        & \multicolumn{2}{c}{All test points} & \multicolumn{2}{c}{Inside} \\
        Degree & Median & Mean & Median & Mean \\
        \midrule
        5  & 0.0276 & 0.0291 & 0.0264 & 0.0278 \\
        10 & 0.0135 & 0.0151 & 0.0126 & 0.0132 \\
        15 & 0.0131 & 0.106  & 0.0123 & 0.0151 \\
        20 & 0.0155 & 19.1   & 0.0136 & 0.0361 \\
        \bottomrule
    \end{tabular}
\end{table}

\subsubsection{Dependence on the number of points and the noise}

Figure~\ref{fig:a_n_sigma} shows the median test MSE up to degree 25. Up to degree 6 the curves for all $n$ coincide: the error is dominated by bias, and more data does not help. The best degree moves up with $n$: 6, 10, 14 and 20 for $n = 30$, 50, 100 and 1000. With $n = 30$ (24 training points) the error grows quickly from about degree 12, since each extra coefficient absorbs more of the noise as the number of parameters approaches the number of training points, and the gaps are larger. With $n = 1000$ the median stays between 0.0102 and 0.0104 from degree 14 to 25: the model has reached the noise floor, and there is no clear optimum.

The best degree falls with the noise: 22, 14 and 6 for $\sigma = 0.01$, 0.1 and 0.5, and for $\sigma = 0.5$ the test error rises from degree 7. Larger $\sigma$ and smaller $n$ act in the same direction. At degree 1 the model is in practice the mean ($x$ is odd and adds nothing), and the three curves start at 0.078, 0.087 and 0.32, differing by the differences in $\sigma^2$. For $\sigma = 0.5$ the test MSE falls from 0.32 to 0.263 at degree 6 against the floor 0.25, so most of the removable error is removed. For $\sigma = 0.01$ the lowest median, $2.0\times10^{-4}$, is twice $\sigma^2$: with little noise the remaining error is mainly bias.

\begin{figure*}
    \centering
    \includegraphics[width=0.49\textwidth]{a_mse_n.pdf}
    \includegraphics[width=0.49\textwidth]{a_mse_sigma.pdf}
    \caption{Median test MSE over 200 repetitions against polynomial degree. Left: $n = 30$, 50, 100 and 1000 with $\sigma = 0.1$ (dotted: $\sigma^2$). Right: $\sigma = 0.01$, 0.1 and 0.5 with $n = 100$, training (dashed) and test (solid), dotted lines at $\sigma^2$. Bands: interquartile range of the test MSE. Open circles: the degree with the lowest median.}
    \label{fig:a_n_sigma}
\end{figure*}

\subsubsection{Coefficients and scaling}\label{sec:res_scaling}

The largest $|\theta_j|$ on the standardised columns grows from 0.72 at degree 5 to 214 at degree 15 and $5\times10^{5}$ at degree 25 (Fig.~\ref{fig:a_theta}). Since all standardised columns have standard deviation 1, this is not caused by small values of $x^j$ but by collinearity: columns such as $x^{13}$ and $x^{15}$ are nearly equal on $[-1,1]$, and OLS gives them large coefficients of opposite sign that nearly cancel. The heat map shows alternating signs, and the coefficients of the low powers also grow with the degree. At degree 15 the largest odd coefficient is 214, but all odd terms together contribute at most 0.15 to the fitted curve. At high degree the individual coefficients cannot be interpreted; only the curve they form has meaning.

\begin{figure*}
    \centering
    \includegraphics[width=0.85\textwidth]{a_theta_degree.pdf}
    \caption{OLS coefficients for seed 42. Left: every coefficient $\theta_j$ on the standardised columns (rows) for each model degree (columns), symmetric logarithmic colour scale. Right: the largest $|\theta_j|$ against degree, on the standardised columns and in the original polynomial; the dotted line marks the range of the left panel.}
    \label{fig:a_theta}
\end{figure*}

Table~\ref{tab:a_scaling} compares scaled and unscaled OLS on the same data. The predictions agree to $10^{-12}$ up to degree 15, since OLS compensates any rescaling of a column by rescaling its coefficient. The condition number is only halved by the scaling but grows by eight orders of magnitude with the degree: on $[-1,1]$ the powers differ little in scale, and what makes $\mathbf{X}$ ill-conditioned is the similarity of the columns, which scaling does not remove. At degree 25 ($\kappa \approx 10^9$) round-off shows as a difference of $2\times10^{-7}$ between the two solutions. For OLS on this data set scaling is therefore not needed. It would matter more on a wider interval such as $[-3, 3]$ or with features in different units. The growing condition number is the problem Ridge addresses.

\begin{table}
    \centering
    \caption{Scaled against unscaled OLS (seed 42). Unscaled: column of ones kept, no centering. $\kappa$: condition number of the training design matrix. $\Delta$: largest difference between the two predictions on the test points.}
    \label{tab:a_scaling}
    \begin{tabular}{rccc}
        \toprule
        & \multicolumn{2}{c}{$\kappa(\mathbf{X})$} & \\
        Degree & Scaled & Unscaled & $\Delta$ \\
        \midrule
        5  & 25                 & 48                 & $2\times10^{-14}$ \\
        10 & $1.7\times10^{3}$  & $3.8\times10^{3}$  & $3\times10^{-14}$ \\
        15 & $1.6\times10^{5}$  & $3.5\times10^{5}$  & $7\times10^{-13}$ \\
        20 & $1.6\times10^{7}$  & $3.6\times10^{7}$  & $1\times10^{-10}$ \\
        25 & $2.7\times10^{9}$  & $6.1\times10^{9}$  & $2\times10^{-7}$  \\
        \bottomrule
    \end{tabular}
\end{table}

\subsection{Ridge regression}\label{sec:res_ridge}

\subsubsection{Shrinkage of the singular-value directions}

Figure~\ref{fig:b_svd} connects Ridge to the singular values of our own design matrix (degree 15, seed 42). The squared singular values fall from 783 to $3\times10^{-8}$ (panel 1). The smallest ones come from collinearity between the high powers and are the flat directions of the cost. Each Ridge penalty $n\lambda$ is a horizontal line in panel 1, and where $\sigma_i^2$ crosses it the shrinkage factor in panel 2 is $1/2$. For $\lambda = 10^{-2}$ ($n\lambda = 0.8$) the crossing lies between directions 6 and 7; the factor is 0.87 for direction 5, 0.33 for direction 7, 0.09 for direction 8, and practically zero from direction 10. The smaller penalties show the same pattern for later directions.

Panel 3 shows how much of the data lies in each pattern $\mathbf{u}_i$. For the first ten directions $|\mathbf{u}_i^T\mathbf{y}|$ is between 0.4 and 1.6, while from direction 11 it lies at the noise level $\sigma = 0.1$: in the flattest directions the data contain practically only noise. Panel 4 shows the consequence. OLS divides this noise by the tiny $\sigma_i$, and the component of $\hat{\boldsymbol{\theta}}$ along $\mathbf{v}_i$ grows from 0.03 to 288; the large components are amplified noise, i.e.\ the variance from Sec.~\ref{sec:res_ols}. Ridge multiplies each component by its shrinkage factor, which we can read off directly: for $\lambda = 10^{-2}$ a third remains in direction 7 and a tenth in direction 8. Ridge thus removes mainly directions that contain noise, which is why it can reduce the variance much at a small cost in bias. For a large penalty it also removes directions with real signal: with $\lambda = 10^{-2}$ directions 6--8 have $|\mathbf{u}_i^T\mathbf{y}|$ between 0.45 and 1.0, well above the noise, and are strongly shrunk.

\begin{figure*}
    \centering
    \includegraphics[width=0.8\textwidth]{b_svd_shrinkage.pdf}
    \caption{Singular-value analysis of the standardised design matrix, degree 15, seed 42, $n_{\mathrm{train}} = 80$. 1: $\sigma_i^2$ with the penalty $n\lambda$ for three values of $\lambda$. 2: the shrinkage factor $\sigma_i^2/(\sigma_i^2+n\lambda)$. 3: $|\mathbf{u}_i^T\mathbf{y}|$, how much of the (centred) training data lies in pattern $\mathbf{u}_i$, with the noise level $\sigma$. 4: the component of $\hat{\boldsymbol{\theta}}$ along $\mathbf{v}_i$ for OLS and Ridge.}
    \label{fig:b_svd}
\end{figure*}

\subsubsection{Error against degree and penalty}

Figures~\ref{fig:b_mse_r2} and \ref{fig:b_heatmap} show the test error for Ridge. With $n = 100$ Ridge does not improve on the best OLS model: the lowest median over all degrees and penalties is OLS at degree 14 (0.0125); $\lambda = 10^{-6}$ is close (0.0133 at degree 10), while $\lambda = 10^{-2}$ gives 0.026. With $\lambda = 10^{-2}$ the shrinkage removes real signal (above), which is bias. What Ridge changes is the behaviour away from the optimum. The Ridge curves do not rise at high degree, and their interquartile bands stay narrow, because the directions that let the curve swing in gaps and at the edges are removed. The means show the same: at degree 20 the mean test MSE is 19 for OLS, 0.037 for $\lambda = 10^{-6}$ and 0.021 for $\lambda = 10^{-4}$. Ridge does not make a good model better, but it makes the result much less sensitive to a poor choice of complexity, which matters because the right degree is not known in practice.

With $n = 30$ the difference is larger (right panel of Fig.~\ref{fig:b_heatmap}). The best median is now a Ridge model, $\lambda \approx 3\times10^{-5}$ at degree 8 (0.0187, against 0.0194 for OLS at degree 6), and OLS explodes from about degree 14, while the rows with small and moderate penalties stay between 0.02 and 0.05. With fewer points there is more noise per parameter, and the penalty pays off more, in line with the dependence on $n$ in Sec.~\ref{sec:res_ols}. The heat map shows medians. The means reveal that a small penalty does not fully protect against the extreme repetitions: at degree 10 the mean test MSE is 3.1 for $\lambda = 10^{-6}$, 0.41 for $\lambda = 10^{-4}$ and 0.044 for $\lambda = 10^{-2}$.

For large penalties the rows of the heat map become nearly uniform. At $\lambda = 10$ and $n = 100$ the median test MSE is 0.073--0.087 at all degrees, close to the degree-1 model, which in practice predicts the mean: the penalty has pushed all coefficients towards zero, and the error is dominated by bias.

\begin{figure*}
    \centering
    \includegraphics[width=0.85\textwidth]{b_mse_r2_degree.pdf}
    \caption{Test MSE (left, logarithmic) and median test $R^2$ (right) against degree for OLS and Ridge with three penalties, $n = 100$, $\sigma = 0.1$, 200 repetitions. Bands: interquartile range. Open circles: the degree with the lowest median. Dotted line: $\sigma^2$.}
    \label{fig:b_mse_r2}
\end{figure*}

\begin{figure*}
    \centering
    \includegraphics[width=0.85\textwidth]{b_heatmap.pdf}
    \caption{Median test MSE over 200 repetitions against degree and penalty, $n = 100$ (left) and $n = 30$ (right), $\sigma = 0.1$. The bottom row is OLS. Colours are clipped at 0.1. Circle: the lowest median.}
    \label{fig:b_heatmap}
\end{figure*}

\subsubsection{Coefficients}

Figure~\ref{fig:b_theta} repeats the coefficient plot of Fig.~\ref{fig:a_theta} for Ridge. While the largest OLS coefficient keeps growing with the degree, the Ridge coefficients level off at about 10, 3.5 and 0.63 for $\lambda = 10^{-6}$, $10^{-4}$ and $10^{-2}$. The penalty lifts the flat directions along which the OLS solution can wander when it fits noise. Once the directions added by a higher degree have $\sigma_i^2 \ll n\lambda$ they are shrunk to nothing, so a larger penalty makes the curve level off at a lower degree and a lower value.

\begin{figure}
    \centering
    \includegraphics[width=\linewidth]{b_theta_degree.pdf}
    \caption{Largest $|\theta_j|$ on the standardised columns against degree for OLS and Ridge (seed 42).}
    \label{fig:b_theta}
\end{figure}

\subsection{Bias--variance trade-off with the bootstrap}\label{sec:res_bootstrap}

Figure~\ref{fig:c_train_test} shows training and test MSE against degree for $n = 50$, in the style of Fig.~2.11 of Hastie \textit{et al.} \cite{Hastie2009}. The training MSE falls monotonically and is below $\sigma^2$ from degree 8. The test MSE has its minimum at degrees 8--10 (0.017) and rises to 0.12 at degree 20 and 39 at degree 25. To the left of the minimum the error is dominated by bias, to the right by variance.

\begin{figure}
    \centering
    \includegraphics[width=\linewidth]{c_train_test.pdf}
    \caption{Training and test MSE against polynomial degree for OLS, $n = 50$, $\sigma = 0.1$. Median over 200 repetitions; bands: interquartile range. Dotted line: $\sigma^2$.}
    \label{fig:c_train_test}
\end{figure}

Figure~\ref{fig:c_bias_variance} shows the bootstrap decomposition for $n = 50$, 100 and 400. For $n = 100$ the squared bias falls from 0.093 at degree 1 to about 0.015 at degree 8 and then stays between 0.014 and 0.018 up to degree 14. The variance is 0.002--0.004 up to degree 8 and then grows by two orders of magnitude, to 0.31 at degree 14. It exceeds the squared bias from degree 11, and the error has its minimum at degree 8 (0.019). With $n = 50$ the variance takes over already at degree 7 and the minimum is at degree 4; the curves are jagged because of the odd and even degrees and the small test set. With $n = 400$ the variance stays below $9\times10^{-4}$ at all degrees, never exceeds the bias, and the error decreases to 0.011 at degree 15. More data thus moves the point where variance takes over to higher degree, in line with Sec.~\ref{sec:res_ols}.

The dashed lines show the squared bias measured against $f$ instead of $y$. For $n = 100$ and 400 the difference between the two is 0.009--0.011 at all degrees where the variance is small, i.e.\ $\sigma^2$: the measured bias absorbs the noise, as expected from Sec.~\ref{sec:biasvar}. The bias measured against $f$ keeps falling with the degree (to $9\times10^{-4}$ at degree 14 for $n = 400$), while the measured bias stops at the noise floor. At high degree and small $n$ the measured bias also grows, because the mean of the 100 bootstrap predictions is itself pulled away by a few extreme fits.

\begin{figure*}
    \centering
    \includegraphics[width=\textwidth]{c_bias_variance.pdf}
    \caption{Bootstrap estimates of the test error, squared bias and variance against degree for OLS, $n = 50$, 100 and 400, $\sigma = 0.1$, 100 bootstrap samples; median over 20 data sets. Solid bias: measured against $y$; dashed: against the true $f$. Dotted line: $\sigma^2$.}
    \label{fig:c_bias_variance}
\end{figure*}

\subsection{Cross-validation}\label{sec:res_cv}

Figure~\ref{fig:d_cv} compares cross-validation with the bootstrap test error and with the MSE on new data for $n = 100$. Cross-validation follows the reference closely, slightly above it (at degree 10: 0.0144 for $k = 5$, 0.0142 for $k = 10$, 0.0133 on new data). It points to degree 11 ($k = 5$) or 14 ($k = 10$), against 12 for the reference; the curves are nearly flat from degree 10 to 16, so the exact minimum is not well determined. The difference between $k = 5$ and $k = 10$ is small; with $k = 10$ each model is trained on 90 instead of 80 points, which gives a slightly lower error at high degree.

The bootstrap agrees with the other two up to degree 7, but then overestimates the error strongly (0.063 at degree 12, against 0.013 on new data) and points to degree 8. The two methods therefore do not point to the same complexity. The reason is that a bootstrap sample drawn with replacement contains on average only about 63\,\% distinct points ($1 - 1/e$), so each model is fitted on about 50 distinct points instead of 80. The variance grows earlier with the degree, as for a smaller $n$ in Sec.~\ref{sec:res_ols}. The bootstrap is well suited to split the error into bias and variance, while cross-validation, which uses all points for both fitting and validation, gives the better estimate of the error for model selection.

\begin{figure}
    \centering
    \includegraphics[width=\linewidth]{d_cv_vs_bootstrap.pdf}
    \caption{OLS: cross-validation MSE ($k = 5$ and 10), bootstrap test error and the MSE on new data against degree, $n = 100$, $\sigma = 0.1$; median over 20 data sets. Open circles: lowest value of each curve.}
    \label{fig:d_cv}
\end{figure}

For Ridge (Fig.~\ref{fig:d_cv_ridge}) the lowest cross-validation error is at degree 11 with $\lambda = 3\times10^{-8}$ (0.0139), practically the same as OLS at the same degree (0.0140). With $n = 100$ the penalty gives no gain at the best degree, as in Sec.~\ref{sec:res_ridge}, but it does at high degree: at degree 15 the best penalty, $\lambda = 3\times10^{-7}$, gives 0.0148 against 0.0178 for OLS. Large penalties give a high error at all degrees (bias).

\begin{figure}
    \centering
    \includegraphics[width=\linewidth]{d_cv_ridge.pdf}
    \caption{Ridge: median 5-fold cross-validation MSE over 20 data sets against degree and $\lambda$, $n = 100$, $\sigma = 0.1$. Colours clipped at 0.1. Circle: the lowest value.}
    \label{fig:d_cv_ridge}
\end{figure}

\subsection{Gradient descent with a fixed learning rate}\label{sec:res_gd}

The analytic gradient and the JAX gradient agree to $2\times10^{-16}$ for both OLS and Ridge, and 2000 gradient-descent iterations with either gradient give coefficients that differ by at most $2\times10^{-15}$. On a larger problem ($10^5 \times 50$) one reverse-mode gradient took 2.6 times as long as one evaluation of the cost (between 2.4 and 3.1 in repeated runs), in line with the statement that the gradient costs a small multiple of the function.

For OLS at degree 6 the Hessian eigenvalues range from $3.4\times10^{-3}$ to 7.63 ($\kappa = 2.3\times10^3$), so $\gamma_{\max} = 2/\mu_{\max} = 0.262$. Figure~\ref{fig:e_gd} shows that gradient descent converges to the closed-form solution for every $\gamma < \gamma_{\max}$ and diverges just above it: at $1.01\,\gamma_{\max}$ the coefficients blow up within about 1100 iterations. The number of iterations to $\lVert\boldsymbol{\theta} - \hat{\boldsymbol{\theta}}\rVert < 10^{-6}$ falls from 336\,841 at $0.05\,\gamma_{\max}$ to 16\,843 at $\gamma^*$, and agrees within 0.5\,\% with the prediction from the slowest mode, $\ln(\lVert\hat{\boldsymbol{\theta}}\rVert/10^{-6}) / (-\ln\max_i|1 - \gamma\mu_i|)$. For Ridge the prediction is 4\,\% too high, since it counts the whole initial distance $\lVert\hat{\boldsymbol{\theta}}\rVert$ as lying along the slowest mode. The slow convergence is caused by the flat direction ($\mu_{\min}$), while the largest usable learning rate is set by the steepest ($\mu_{\max}$). Ridge with $\lambda = 10^{-2}$ raises $\mu_{\min}$ to 0.023 ($\kappa = 327$) and needs about eight times fewer iterations (2138 at $0.99\,\gamma_{\max}$). The converged solutions reproduce the closed-form test MSE to six digits (0.015751 for OLS, 0.024435 for Ridge).

\begin{figure*}
    \centering
    \includegraphics[width=0.85\textwidth]{e_gd.pdf}
    \caption{Plain gradient descent, degree 6, seed 42. Left: iterations to reach $\lVert\boldsymbol{\theta} - \hat{\boldsymbol{\theta}}\rVert < 10^{-6}$ against the learning rate relative to $\gamma_{\max} = 2/\mu_{\max}$, measured (points) and predicted from the Hessian eigenvalues (lines), for OLS and Ridge ($\lambda = 10^{-2}$). Right: distance to the closed-form solution against iteration.}
    \label{fig:e_gd}
\end{figure*}

\subsection{Momentum and adaptive learning rates}\label{sec:res_adaptive}

Figure~\ref{fig:f_methods} shows the number of iterations needed to bring the cost within $10^{-8}$ of the closed-form minimum, against the (initial) learning rate, for the same problem as in Sec.~\ref{sec:res_gd}. The criterion is now on the cost rather than on the distance to $\hat{\boldsymbol{\theta}}$ as in Sec.~\ref{sec:res_gd}; in a flat direction a large distance costs little, and the cost is what matters for the predictions. For OLS plain gradient descent converges within $10^5$ iterations for $\gamma$ from 0.032 to 0.178, the largest value in our grid below $\gamma_{\max} = 0.262$, and the number of iterations falls as $1/\gamma$, from 66\,559 to 11\,834. Momentum is the fastest method: 601 iterations at best, 20 times fewer than plain gradient descent, and it converges for $\gamma$ from 0.0032 to 0.32. AdaGrad converges for $\gamma \geq 0.018$ but needs at least 13\,786 iterations, and below $\gamma = 0.1$ the number rises steeply (92\,163 at $\gamma = 0.018$), since the sum of squared gradients only grows and the steps keep shrinking. Adam needs 842 iterations at best, but its main advantage is robustness: it reaches the criterion for all 17 learning rates from $10^{-4}$ to 1, and from $\gamma \approx 0.06$ the number of iterations hardly changes (842--1135). Since it divides by the running size of the gradient, a large gradient does not give a large step, and a too large $\gamma$ does not make it diverge as plain gradient descent does. With Ridge ($\lambda = 10^{-2}$) all methods that converge are faster, as in Sec.~\ref{sec:res_gd}, since the flattest direction is lifted: 138 iterations for momentum, 174 for Adam and 1475 for plain gradient descent.

The adaptive methods rescale each coefficient separately, but the flat direction of this problem is a combination of nearly collinear columns, not a coordinate axis (Sec.~\ref{sec:adaptive}). They therefore do not remove the ill-conditioning; momentum, which builds up speed along the valley floor, gives the largest gain.

RMSProp never reaches the criterion. Every coefficient ends in a cycle, jumping between two points $\gamma/2$ on either side of the optimum: when the size of the gradient is constant, the running mean $r$ equals $g^2$, and the step $\gamma g/\sqrt{r}$ is exactly $\gamma$, however small the gradient is. The remaining excess cost therefore grows as $\gamma^2$: $5.7\times10^{-8}$, $5.7\times10^{-6}$ and $5.7\times10^{-4}$ for $\gamma = 10^{-4}$, $10^{-3}$ and $10^{-2}$. Adam reaches the criterion but does not stay there: with $\gamma = 0.056$ the excess cost repeatedly jumps to between $3.8\times10^{-3}$ and $5.3\times10^{-3}$ from iteration 1465 on (right panel). Adam divides the step by a running average of the size of earlier gradients. Near the minimum the gradients are small, this average shrinks slowly, and the steps therefore become gradually longer. In the end they become too long, as for plain gradient descent with $\gamma$ above $\gamma_{\max}$, and the deviation grows quickly until the larger gradients make the steps shorter again. We checked that the steps pass the stability limit shortly before the first jump, at iteration 1366. A smaller $\gamma$ gives smaller jumps ($\gamma = 10^{-3}$ does not exceed $1.3\times10^{-6}$ in the last 1000 iterations), but needs more iterations.

\begin{figure*}
    \centering
    \includegraphics[width=\textwidth]{f_methods.pdf}
    \caption{Degree 6, seed 42, start at $\boldsymbol{\theta} = 0$. Left and middle: iterations until the cost is within $10^{-8}$ of the closed-form minimum against the (initial) learning rate, for OLS and Ridge ($\lambda = 10^{-2}$); missing points did not reach it within $10^5$ iterations (RMSProp never does). Right: excess cost against iteration for OLS at each method's best learning rate (RMSProp at $10^{-3}$), first $3\times10^4$ iterations.}
    \label{fig:f_methods}
\end{figure*}

\subsection{Lasso}\label{sec:res_lasso}

JAX returns $\mathrm{d}|\theta|/\mathrm{d}\theta = 1$ at $\theta = 0$ (also at $-0$), a valid subgradient. For $\lambda = 10^{-4}$ and $10^{-3}$ no coefficient is zero at the optimum, and both plain gradient descent and Adam reproduce scikit-learn's solution (cost difference below $10^{-8}$). For $\lambda = 10^{-2}$ scikit-learn sets $\theta_5$ and $\theta_6$ exactly to zero, and the optimality condition holds ($|\frac{2}{n}\mathbf{x}_j^T\mathbf{r}| = 4.8\times10^{-3}$ and $3.9\times10^{-3}$, below $\lambda$). Gradient descent never produces exact zeros. Where a coefficient is zero at the optimum, the cost has a kink, and the slope does not become smaller close to the kink; with a fixed $\gamma$ the steps therefore do not shrink, and gradient descent keeps jumping across it. Adam comes to within about $1.5\times10^{-6}$ of the optimal cost (smallest $|\theta_j| = 3\times10^{-5}$). Plain gradient descent at $0.99\,\gamma_{\max}$, where it is barely stable, fails: the excess cost stays at 0.12--0.13, and four of the six coefficients changed sign in the last step. A smaller step gives smaller jumps (Table~\ref{tab:g_lasso}), but a fixed-step subgradient method never settles exactly; a decreasing step, or a method that sets coefficients exactly to zero, as scikit-learn's coordinate descent does, is needed.

\begin{table}
    \centering
    \caption{Plain subgradient descent for Lasso (analytic subgradient, $\mathrm{sign}(0) = 0$), $\lambda = 10^{-2}$, degree 6: smallest and largest excess cost over scikit-learn's solution in the last 1000 of $10^5$ iterations, for four fixed learning rates.}
    \label{tab:g_lasso}
    \begin{tabular}{cc}
        \toprule
        $\gamma / \gamma_{\max}$ & Excess cost \\
        \midrule
        0.99 & 0.12--0.13 \\
        0.90 & $5.8\times10^{-4}$--$1.6\times10^{-3}$ \\
        0.50 & $2.1\times10^{-6}$--$3.4\times10^{-5}$ \\
        0.10 & $8.3\times10^{-8}$--$4.5\times10^{-6}$ \\
        \bottomrule
    \end{tabular}
\end{table}

Figure~\ref{fig:g_lasso} shows the Lasso path from scikit-learn with the gradient-descent solutions. As $\lambda$ grows the coefficients shrink and more of them become exactly zero: none at $\lambda = 10^{-3}$, $\theta_5$ and $\theta_6$ at $10^{-2}$, and at $0.1$ only the $x^2$ term is left. The odd powers are small for all $\lambda$, as expected for an even function. On the test set (seed 42) Lasso with $\lambda = 10^{-4}$ and $10^{-3}$ is nearly as good as OLS (0.0158 and 0.0162 against 0.0158); with $\lambda = 10^{-2}$ the error rises to 0.024, like Ridge with the same $\lambda$ (0.024). At degree 6 there is little variance to remove, so neither penalty helps here.

\begin{figure}
    \centering
    \includegraphics[width=\linewidth]{g_lasso_path.pdf}
    \caption{Lasso coefficients on the standardised columns against $\lambda$, degree 6, seed 42. Lines: scikit-learn (blue: even powers, orange: odd powers). Markers: gradient descent with the JAX subgradient at $\lambda = 10^{-4}$, $10^{-3}$ and $10^{-2}$, plain with $\gamma = 0.99\,\gamma_{\max}$ (circles) and Adam with $\gamma = 10^{-3}$ (crosses); at $\lambda = 10^{-2}$ plain gradient descent misses the path.}
    \label{fig:g_lasso}
\end{figure}

\subsection{Stochastic gradient descent}\label{sec:res_sgd}

Figure~\ref{fig:h_sgd} shows the training cost above the closed-form minimum against epochs. With plain SGD and a constant $\gamma = 0.05$, smaller batches, down to $M = 5$, give more updates per epoch and faster progress, and $M = 5$ reaches the lowest cost after 1000 epochs ($1.7\times10^{-4}$, test MSE 0.0164 against 0.0158 for the closed form). $M = 1$ is too noisy: $\gamma = 0.05$ is above the stability limit $2/(2\lVert\mathbf{x}_i\rVert^2)_{\max} = 0.018$ for single points, and the cost fluctuates at about $10^{-2}$ or higher ($1.2\times10^{-2}$ after 100 and $1.8\times10^{-2}$ after 1000 epochs; test MSE 0.058). With $M = 80$ plain SGD is full-batch gradient descent with a small step and is slow. We therefore used $M = 5$ for the comparison of the update rules. The schedule $\gamma_t = 50/(t + 500)$ decays too fast when $t$ counts updates (16 per epoch): after 1000 epochs $\gamma_t = 0.003$, and the cost is $4.6\times10^{-3}$, against $1.7\times10^{-4}$ with the constant $\gamma = 0.05$.

Measured in passes through the data, SGD reaches a moderate accuracy about as fast as full-batch gradient descent with the same update rule, or faster. An excess cost of $10^{-3}$ takes 269 epochs with SGD and momentum and 333 with SGD and Adam, against 397 and 315 iterations for full-batch momentum and Adam, and plain SGD gets there after 523 epochs, while full-batch plain gradient descent is still at $2.6\times10^{-3}$ after 1000. The learning rates were chosen separately, since SGD needs smaller ones, and for SGD the count is the first epoch below the target, which favours a noisy method. But SGD with momentum and Adam does not settle lower: the gradient noise keeps it fluctuating ($2.8\times10^{-3}$ and $1.4\times10^{-3}$ after 1000 epochs for momentum and Adam), while full-batch momentum reaches $1.5\times10^{-5}$ and full-batch Adam $2\times10^{-10}$. With 80 training points and six parameters a full gradient is cheap, and the advantage of SGD, many cheap updates when $n$ is large, does not come into play. SGD with momentum and Adam ended with a slightly lower test MSE on our test set (0.0146 and 0.0149) than the closed form (0.0158) without being converged.

\begin{figure*}
    \centering
    \includegraphics[width=0.85\textwidth]{h_sgd.pdf}
    \caption{SGD for OLS, degree 6, seed 42: training cost above the closed-form minimum against epochs (running median over 25 epochs for SGD). Left: plain SGD with $\gamma = 0.05$ for batch sizes $M$ ($M = 80$ is full batch). Right: SGD with $M = 5$ (constant rate, decaying schedule, Adam) against full-batch gradient descent, per pass through the data.}
    \label{fig:h_sgd}
\end{figure*}

\subsection{Final model selection}\label{sec:res_final}

Table~\ref{tab:i_selected} summarises the models chosen by cross-validation (Sec.~\ref{sec:resampling}), and Fig.~\ref{fig:i_cv} shows the cross-validation error against degree with the best $\lambda$ at each degree. By the cross-validated error the best model is Ridge with a small penalty: median CV MSE 0.0149 at a median degree of 11.5, with $\lambda$ between $10^{-8}$ and $10^{-5}$, against 0.0158 for OLS at degree 10. On new data the methods are nearly equal, 0.0136--0.0141, or about $1.4\,\sigma^2$, and which method is best varies between the data sets (OLS in 6, Ridge in 7 and Lasso in 7 of 20). The differences between the medians, at most 0.0005, are small compared with the spread over data sets: the interquartile range of the MSE on new data is 0.0131--0.0146 for Ridge and 0.0132--0.0150 for OLS. For OLS, Ridge and Lasso the selected degree ranges from 8 to 16. From degree 8 on, the median cross-validation error of Ridge and Lasso stays between 0.0154 and 0.0186, and that of OLS between 0.017 and 0.021 up to degree 12 (Fig.~\ref{fig:i_cv}), so the minimum is poorly determined, as in Sec.~\ref{sec:res_cv}.

\begin{table}
    \centering
    \caption{Models selected by 5-fold cross-validation on the training part of each of 20 data sets ($n = 100$, $\sigma = 0.1$): median and range of the degree, median $\lambda$, and median CV MSE, test MSE and MSE on new data. Lasso$^*$: only $\lambda \geq 10^{-3}$ allowed.}
    \label{tab:i_selected}
    \begin{tabular}{lccccc}
        \toprule
        Method & Degree & $\lambda$ & CV & Test & New data \\
        \midrule
        OLS         & 10 (8--16)   & --                 & 0.0158 & 0.0155 & 0.0141 \\
        Ridge       & 11.5 (8--16) & $10^{-7}$          & 0.0149 & 0.0133 & 0.0136 \\
        Lasso       & 12 (8--16)   & $5.5\times10^{-6}$ & 0.0163 & 0.0139 & 0.0140 \\
        Lasso$^*$   & 11.5 (6--16) & $10^{-3}$          & 0.0207 & 0.0184 & 0.0185 \\
        \bottomrule
    \end{tabular}
\end{table}

\begin{figure}
    \centering
    \includegraphics[width=\linewidth]{i_cv_methods.pdf}
    \caption{Median 5-fold CV MSE over 20 data sets against degree, with the best $\lambda$ at each degree for Ridge and Lasso; bands: interquartile range. For Lasso with $\lambda < 10^{-3}$ the solver has not converged at high degree (see text). $n = 100$, $\sigma = 0.1$.}
    \label{fig:i_cv}
\end{figure}

The methods differ at high degree. Up to degree 9 the curves of OLS, Ridge and Lasso differ by at most 12\,\%; above degree 11 the OLS error rises, to 0.036 at degree 16 with a wide band, while Ridge stays between 0.0154 and 0.0175 from degree 10 on. The strength of Ridge for this problem is that a small penalty removes the unstable directions (Sec.~\ref{sec:res_ridge}) and makes the result insensitive to the choice of degree, at practically no cost, and the solution has a closed form. OLS needs no $\lambda$, but its error depends strongly on getting the degree right.

For Lasso, cross-validation chose the smallest penalties in the grid ($10^{-6}$ in 10 of 20 data sets), and only 4 of the 20 selected models have a coefficient set to zero. At these penalties scikit-learn's coordinate descent often did not converge within $10^5$ iterations: 18 of the 20 selected Lasso fits stopped at the iteration limit, with a median duality gap of $8\times10^{-3}$ against a cost of $1.3\times10^{-2}$. What we call Lasso there is an early-stopped approximation, which regularises in its own way (Sec.~\ref{sec:res_sgd}), and its similarity to Ridge says little about the Lasso solution itself. If only $\lambda \geq 10^{-3}$ is allowed, where the solver converges (duality gap $2\times10^{-9}$), cross-validation chooses $\lambda = 10^{-3}$ in all 20 data sets. The models then have a median of 6 nonzero coefficients (16 of 20 have at least one zero), but the error on new data rises to 0.0185. For this problem the sparsity of Lasso is thus not an advantage: a good polynomial approximation of Runge's function needs many terms, and a penalty large enough to remove terms also adds bias. Its weaknesses here are the lack of a closed form and a slow solver when the columns are collinear.

\begin{table}
    \centering
    \caption{Bootstrap decomposition at degree 14 with the median selected penalties, 100 bootstrap samples, median over 20 data sets. Bias$^2$ is measured against the true $f$.}
    \label{tab:i_bootstrap}
    \begin{tabular}{lccc}
        \toprule
        Method & Error & Bias$^2$ & Variance \\
        \midrule
        OLS                                     & 0.102  & $8.0\times10^{-3}$ & 0.087 \\
        Ridge, $\lambda = 10^{-7}$              & 0.036  & $3.1\times10^{-3}$ & 0.024 \\
        Lasso, $\lambda = 5.5\times10^{-6}$     & 0.024  & $4.0\times10^{-3}$ & 0.013 \\
        Lasso, $\lambda = 10^{-3}$              & 0.021  & $8.9\times10^{-3}$ & 0.0023 \\
        \bottomrule
    \end{tabular}
\end{table}

In terms of the bias--variance decomposition, a higher degree in OLS attacks the bias: the bias against $f$ falls with the degree (Sec.~\ref{sec:res_bootstrap}). Ridge and Lasso attack the variance. At degree 14 the variance falls from 0.087 for OLS to 0.024 with Ridge ($\lambda = 10^{-7}$) and to 0.0023 with the converged Lasso ($\lambda = 10^{-3}$), see Table~\ref{tab:i_bootstrap}. For the converged Lasso the price is visible in the bias, $8.9\times10^{-3}$ against $3.1\times10^{-3}$ for Ridge. The bias of OLS, $8.0\times10^{-3}$, is not a real bias but the effect of a few extreme bootstrap fits on the mean prediction (Sec.~\ref{sec:res_bootstrap}). The numbers thus support the reading that the penalties trade variance for bias, with $\lambda$ setting the balance. At degree 14 the bootstrap ranks the converged Lasso best, unlike cross-validation; since the bootstrap overestimates the variance (Sec.~\ref{sec:res_cv}), it favours strong penalties.
